%% ╔══════════════════════════════════════════════════════════════════╗
%% ║         CORIOTLAB — PLANTILLA INFORME TÉCNICO EXTERNO           ║
%% ║         Nivel consultoría profesional                           ║
%% ╚══════════════════════════════════════════════════════════════════╝
%%
%% COMPILAR: xelatex -interaction=nonstopmode plantilla_tecnico.tex
%%           (doble pasada — ejecutar el comando dos veces)

%% ═══════════════════════════════════════════════════════════════════
%% ZONA DE CONFIGURACIÓN — editar aquí antes de compilar
%% ═══════════════════════════════════════════════════════════════════

%% TÍTULO completo del informe técnico (máx. 100 caracteres)
\newcommand{\DocTitulo}{[Título completo del informe técnico]}

%% SUBTÍTULO — enfoque metodológico o alcance
\newcommand{\DocSubtitulo}{[Subtítulo técnico: enfoque o alcance del proyecto]}

%% NÚMERO DE DOCUMENTO — Formato: IT-AAAA-NNN
\newcommand{\DocNumero}{IT-2025-[NNN]}

%% PROYECTO — nombre oficial
\newcommand{\DocProyecto}{[Nombre oficial del proyecto]}

%% CÓDIGO INTERNO — Formato: CORIOT-AAAA-SIGLAS-NNN
\newcommand{\DocCodigo}{CORIOT-2025-[SIGLAS]-[NNN]}

%% CLIENTE / ENTIDAD destinataria
\newcommand{\DocCliente}{[Nombre del cliente o entidad destinataria]}

%% PERÍODO — rango de fechas del proyecto
\newcommand{\DocPeriodo}{[Mes inicial] --- [Mes final] [AAAA]}

%% LUGAR
\newcommand{\DocLugar}{Medellín, Colombia}

%% VERSIÓN — iniciar en 1.0
\newcommand{\DocVersion}{1.0}

%% FECHA de emisión — DD de mes de AAAA
\newcommand{\DocFecha}{[DD de mes de AAAA]}

%% CONFIDENCIALIDAD: "Uso interno --- CORIOTLAB" | "Confidencial" | "Público"
\newcommand{\DocConfidencialidad}{Uso interno --- CORIOTLAB}

%% AUTOR principal
\newcommand{\DocAutor}{[Nombre completo del autor]}

%% CARGO DEL AUTOR
\newcommand{\DocCargoAutor}{[Cargo --- CORIOTLAB]}

%% REVISOR
\newcommand{\DocRevisor}{[Nombre del revisor]}

%% CARGO DEL REVISOR
\newcommand{\DocCargoRevisor}{[Cargo del revisor --- CORIOTLAB]}

%% ENTIDAD
\newcommand{\DocEntidad}{CORIOTLAB --- ITM}

%% ═══════════════════════════════════════════════════════════════════

\documentclass[12pt, letterpaper]{article}

%% --- Codificación y lenguaje
\usepackage[spanish, es-nodecimaldot, es-noshorthands]{babel}

%% --- Fuentes (XeLaTeX — por nombre de familia Windows)
\usepackage{fontspec}
\setmainfont{Inter}
\newfontfamily\FuenteTitulo{MuseoModerno}
\newfontfamily\FuenteCodigo{Space Mono}

%% --- Geometría — márgenes ajustados al membrete (header ~3cm, footer ~2.5cm)
\usepackage[
  top=3.5cm, bottom=3.0cm,
  left=2.5cm, right=2.5cm
]{geometry}

%% --- Color
\usepackage[table]{xcolor}
\definecolor{AzulITM}     {HTML}{102D69}
\definecolor{Magenta}     {HTML}{C14894}
\definecolor{AzulDigital} {HTML}{56ACDE}
\definecolor{GrisPizarra} {HTML}{2F2F2F}
\definecolor{GrisClaro}   {HTML}{F2F2F2}
\definecolor{GrisMedio}   {HTML}{AAAAAA}
\definecolor{GrisLinea}   {HTML}{DDDDDD}
\definecolor{Blanco}      {HTML}{FFFFFF}

%% --- Gráficos
\usepackage{graphicx}
\usepackage{tikz}
\usepackage{eso-pic}

%% --- Tablas
\usepackage{longtable}
\usepackage{booktabs}
\usepackage{array}
\usepackage{colortbl}
\usepackage{multirow}

%% --- Listas
\usepackage{enumitem}

%% --- Cajas
\usepackage[most]{tcolorbox}

%% --- Código
\usepackage{listings}

%% --- TOC y secciones
\usepackage{titlesec}
\usepackage{titletoc}
\usepackage[hidelinks]{hyperref}

%% --- Espaciado
\usepackage{parskip}
\setlength{\parskip}{6pt}
\setlength{\parindent}{0pt}
\renewcommand{\arraystretch}{1.4}

%% -----------------------------------------------------------------
%% MEMBRETE — opacidad completa en todas las páginas
%% -----------------------------------------------------------------
\AddToShipoutPictureBG{%
  \begin{tikzpicture}[remember picture, overlay]
    \node at (current page.center) {%
      \includegraphics[width=\paperwidth, height=\paperheight,
        keepaspectratio=false]%
        {../../membretes/membrete_azul.png}%
    };
  \end{tikzpicture}%
}

%% Sin encabezado ni pie — el membrete los provee
\pagestyle{empty}

%% -----------------------------------------------------------------
%% SECCIONES
%% -----------------------------------------------------------------
\titleformat{\section}
  {\FuenteTitulo\large\color{AzulITM}\bfseries}
  {\thesection.}{0.5em}{}
\titleformat{\subsection}
  {\FuenteTitulo\normalsize\color{GrisPizarra}\bfseries}
  {\thesubsection.}{0.5em}{}
\titleformat{\subsubsection}
  {\normalsize\color{GrisPizarra}\bfseries\itshape}
  {\thesubsubsection.}{0.5em}{}
\titlespacing*{\section}      {0pt}{2em}{0.8em}
\titlespacing*{\subsection}   {0pt}{1.2em}{0.4em}
\titlespacing*{\subsubsection}{0pt}{0.8em}{0.3em}

%% -----------------------------------------------------------------
%% CAJAS tcolorbox
%% -----------------------------------------------------------------
\tcbset{
  cajaNota/.style={
    colback=AzulDigital!10, colframe=AzulDigital,
    fonttitle=\FuenteTitulo\small\bfseries, coltitle=white,
    boxrule=0pt, leftrule=4pt, title=Nota,
    left=8pt, top=6pt, bottom=6pt
  },
  cajaAlerta/.style={
    colback=Magenta!8, colframe=Magenta,
    fonttitle=\FuenteTitulo\small\bfseries, coltitle=white,
    boxrule=0pt, leftrule=4pt, title=Advertencia,
    left=8pt, top=6pt, bottom=6pt
  },
  cajaResultado/.style={
    colback=AzulITM!8, colframe=AzulITM,
    fonttitle=\FuenteTitulo\small\bfseries, coltitle=white,
    boxrule=0pt, leftrule=4pt, title=Resultado,
    left=8pt, top=6pt, bottom=6pt
  },
  cajaDatos/.style={
    colback=GrisClaro, colframe=GrisLinea,
    boxrule=0.6pt, left=8pt, top=6pt, bottom=6pt
  },
  cajaResumen/.style={
    colback=AzulITM!8, colframe=AzulITM,
    fonttitle=\FuenteTitulo\normalsize\bfseries, coltitle=white,
    boxrule=1pt, title=Resumen Ejecutivo,
    left=10pt, right=10pt, top=8pt, bottom=8pt
  }
}

%% -----------------------------------------------------------------
%% CÓDIGO lstlisting
%% -----------------------------------------------------------------
\lstdefinestyle{coriotcode}{
  basicstyle=\FuenteCodigo\small,
  backgroundcolor=\color{GrisClaro},
  commentstyle=\color{GrisMedio}\itshape,
  keywordstyle=\color{AzulITM}\bfseries,
  stringstyle=\color{Magenta},
  numberstyle=\tiny\color{GrisMedio},
  numbers=left, numbersep=10pt,
  breaklines=true, frame=none,
  tabsize=2, showstringspaces=false,
  aboveskip=8pt, belowskip=8pt
}
\lstset{style=coriotcode}

%% =================================================================
%%  INICIO DEL DOCUMENTO
%% =================================================================
\begin{document}
\color{GrisPizarra}

%% -----------------------------------------------------------------
%% 1. PORTADA — misma lógica que el Word: título + metadata, sin barra TikZ
%% -----------------------------------------------------------------

{\FuenteTitulo\LARGE\color{AzulITM}\bfseries\DocTitulo}\par\vspace{0.3em}
{\FuenteTitulo\normalsize\color{GrisPizarra}\DocSubtitulo}\par\vspace{1em}
{\color{GrisLinea}\rule{\linewidth}{0.5pt}}\par\vspace{0.8em}
\begin{tabular}{@{}ll@{}}
  {\small\color{GrisMedio}Proyecto:}          & {\small\color{GrisPizarra}\DocProyecto} \\[3pt]
  {\small\color{GrisMedio}Código:}            & {\small\color{GrisPizarra}\DocCodigo} \\[3pt]
  {\small\color{GrisMedio}Cliente / Entidad:} & {\small\color{GrisPizarra}\DocCliente} \\[3pt]
  {\small\color{GrisMedio}Período:}           & {\small\color{GrisPizarra}\DocPeriodo} \\[3pt]
  {\small\color{GrisMedio}Lugar:}             & {\small\color{GrisPizarra}\DocLugar} \\[3pt]
  {\small\color{GrisMedio}Versión:}           & {\small\color{GrisPizarra}\DocVersion} \\[3pt]
  {\small\color{GrisMedio}Fecha:}             & {\small\color{GrisPizarra}\DocFecha} \\[3pt]
  {\small\color{GrisMedio}Confidencialidad:}  & {\small\color{GrisPizarra}\DocConfidencialidad} \\
\end{tabular}

\vspace{1.4em}

%% RESUMEN EJECUTIVO en portada
\begin{tcolorbox}[cajaResumen]
  {[}Síntesis ejecutiva del informe. Describa en 5--8 oraciones:
  el contexto del proyecto, el problema abordado, la metodología
  empleada, los principales resultados obtenidos y la conclusión
  central. Evite tecnicismos excesivos: este resumen lo leerán
  directivos y clientes sin formación técnica profunda.{]}
\end{tcolorbox}

\newpage

%% -----------------------------------------------------------------
%% 2. CONTROL DE VERSIONES
%% -----------------------------------------------------------------
\section*{Control de Versiones}
\addcontentsline{toc}{section}{Control de Versiones}

\setlength{\LTpre}{4pt}\setlength{\LTpost}{4pt}
\begin{longtable}{
  >{\centering\arraybackslash}p{1.4cm}
  >{\centering\arraybackslash}p{2.6cm}
  p{3.8cm}
  p{6.4cm}
}
\rowcolor{AzulITM}
\color{Blanco}\FuenteTitulo\small\bfseries Versión &
\color{Blanco}\FuenteTitulo\small\bfseries Fecha &
\color{Blanco}\FuenteTitulo\small\bfseries Autor &
\color{Blanco}\FuenteTitulo\small\bfseries Descripción del cambio \\
\endfirsthead
\rowcolor{AzulITM}
\color{Blanco}\FuenteTitulo\small\bfseries Versión &
\color{Blanco}\FuenteTitulo\small\bfseries Fecha &
\color{Blanco}\FuenteTitulo\small\bfseries Autor &
\color{Blanco}\FuenteTitulo\small\bfseries Descripción del cambio \\
\endhead
\rowcolor{white}
1.0 & {[}DD/MM/AAAA{]} & \DocAutor & Versión inicial. \\
\rowcolor{GrisClaro}
1.1 & {[}DD/MM/AAAA{]} & \DocAutor & {[}Descripción de la revisión.{]} \\
\end{longtable}

%% -----------------------------------------------------------------
%% 3. TABLA DE CONTENIDO
%% -----------------------------------------------------------------
\newpage
\tableofcontents
\newpage

%% -----------------------------------------------------------------
%% 4. INTRODUCCIÓN
%% -----------------------------------------------------------------
\section{Introducción}

{\itshape\color{GrisMedio}%
{[}Párrafo 1: contexto institucional del proyecto. ¿En qué marco se
desarrolla este trabajo? ¿Qué área o entidad lo impulsa?{]}

{[}Párrafo 2: descripción del problema o necesidad que motivó el proyecto.
¿Qué situación se busca resolver o mejorar?{]}

{[}Párrafo 3: descripción general de la solución implementada y el
enfoque técnico adoptado.{]}}

\begin{tcolorbox}[cajaDatos]
  \textbf{Alcance del documento:} {[}Especifique qué fases, sistemas
  o componentes cubre este informe, y cuáles quedan fuera del alcance.
  Indique el período que comprende: \DocPeriodo.{]}
\end{tcolorbox}

%% -----------------------------------------------------------------
%% 5. OBJETIVOS
%% -----------------------------------------------------------------
\section{Objetivos}

\subsection{Objetivo General}

{\itshape\color{GrisMedio}%
{[}Verbo infinitivo{]} {[}objeto o sistema desarrollado{]} para {[}propósito
o beneficio esperado{]}, mediante {[}método o tecnología empleada{]}.}

\subsection{Objetivos Específicos}

\begin{itemize}[label={\color{AzulITM}\textbf{>}}, leftmargin=1.8em, itemsep=5pt]
  \item {[}Objetivo específico 1: acción concreta y verificable.{]}
  \item {[}Objetivo específico 2: acción concreta y verificable.{]}
  \item {[}Objetivo específico 3: acción concreta y verificable.{]}
\end{itemize}

%% -----------------------------------------------------------------
%% SECCIONES OPCIONALES — descomentar según proyecto
%% -----------------------------------------------------------------

%% ── BLOQUE OPT-A: Marco Teórico ──────────────────────────────────
%\section{Marco Teórico}
%\subsection{{[}Concepto o tecnología clave 1{]}}
%{[}Definición y relevancia para el proyecto. 1--2 párrafos.{]}
%\subsection{{[}Concepto o tecnología clave 2{]}}
%{[}Definición y relevancia para el proyecto. 1--2 párrafos.{]}

%% ── BLOQUE OPT-B: Metodología ────────────────────────────────────
%\section{Metodología}
%{[}Descripción del enfoque metodológico adoptado y justificación.{]}
%\begin{enumerate}[leftmargin=1.8em, itemsep=4pt]
%  \item \textbf{Fase 1 --- {[}Nombre{]}:} {[}Descripción de actividades.{]}
%  \item \textbf{Fase 2 --- {[}Nombre{]}:} {[}Descripción de actividades.{]}
%  \item \textbf{Fase 3 --- {[}Nombre{]}:} {[}Descripción de actividades.{]}
%\end{enumerate}

%% ── BLOQUE OPT-C: Desarrollo Técnico + Código ────────────────────
%\section{Desarrollo Técnico}
%\subsection{{[}Nombre del componente 1{]}}
%{[}Descripción del componente: propósito, diseño e implementación.{]}
%\begin{tcolorbox}[cajaDatos]
%\begin{lstlisting}[language=Python]
%def ejemplo():
%    pass
%\end{lstlisting}
%\end{tcolorbox}
%\subsection{{[}Nombre del componente 2{]}}
%{[}Descripción del componente.{]}
%\begin{tcolorbox}[cajaDatos]
%\begin{lstlisting}[language=C++]
%void setup() { }
%\end{lstlisting}
%\end{tcolorbox}

%% ── BLOQUE OPT-D: Resultados con tabla de indicadores ────────────
%\section{Resultados}
%\begin{longtable}{p{4.2cm} >{\centering}p{2.2cm} >{\centering}p{2.2cm}
%    >{\centering}p{2.2cm} p{3.2cm}}
%\rowcolor{AzulITM}
%\color{Blanco}\FuenteTitulo\small\bfseries Indicador &
%\color{Blanco}\FuenteTitulo\small\bfseries Meta &
%\color{Blanco}\FuenteTitulo\small\bfseries Logrado &
%\color{Blanco}\FuenteTitulo\small\bfseries \% Cumpl. &
%\color{Blanco}\FuenteTitulo\small\bfseries Observaciones \\
%\endfirsthead
%\rowcolor{white}
%{[}Indicador 1{]} & {[}Meta{]} & {[}Valor{]} & {[}XX\%{]} & {[}Nota.{]} \\
%\rowcolor{GrisClaro}
%{[}Indicador 2{]} & {[}Meta{]} & {[}Valor{]} & {[}XX\%{]} & {[}Nota.{]} \\
%\end{longtable}
%\begin{tcolorbox}[cajaResultado]
%  {[}Síntesis de los resultados.{]}
%\end{tcolorbox}

%% ── BLOQUE OPT-E: Recomendaciones ───────────────────────────────
%\section{Recomendaciones}
%\begin{itemize}[label={\color{AzulITM}\textbf{>}}, leftmargin=1.8em, itemsep=4pt]
%  \item {[}Recomendación 1: acción concreta, justificación breve.{]}
%  \item {[}Recomendación 2: acción concreta, justificación breve.{]}
%\end{itemize}

%% ── BLOQUE OPT-F: Referencias Bibliográficas ─────────────────────
%\section{Referencias Bibliográficas}
%\begin{enumerate}[leftmargin=1.8em, itemsep=4pt]
%  \item {[}Apellido, Inicial. (Año). \textit{Título}. Editorial / URL.{]}
%\end{enumerate}

%% ── BLOQUE OPT-G: Anexos ─────────────────────────────────────────
%\appendix
%\section{{[}Nombre del Anexo A{]}}
%{[}Contenido del anexo A.{]}

%% -----------------------------------------------------------------
%% 6. CONCLUSIONES
%% -----------------------------------------------------------------
\section{Conclusiones}

{\itshape\color{GrisMedio}%
{[}Párrafo 1: valoración global del proyecto. ¿Se cumplieron los objetivos?{]}

{[}Párrafo 2: hallazgos técnicos, limitaciones o aprendizajes relevantes.{]}}

\begin{tcolorbox}[cajaResultado]
  {[}Síntesis evaluativa: el proyecto {[}cumplió / superó / no alcanzó{]}
  los indicadores de desempeño establecidos.{]}
\end{tcolorbox}

%% -----------------------------------------------------------------
%% 7. FIRMAS Y APROBACIÓN
%% -----------------------------------------------------------------
\section*{Firmas y Aprobación}
\addcontentsline{toc}{section}{Firmas y Aprobación}

\vspace{1.5em}
\noindent
\begin{tabular}{p{7cm} p{1cm} p{7cm}}
  \multicolumn{1}{c}{\textbf{\small\color{AzulITM}Elaborado por}} & &
  \multicolumn{1}{c}{\textbf{\small\color{AzulITM}Revisado y aprobado por}} \\[0.6em]
  \multicolumn{1}{l}{\color{AzulITM}\rule{7cm}{0.6pt}} & &
  \multicolumn{1}{l}{\color{AzulITM}\rule{7cm}{0.6pt}} \\[4pt]
  {\small\bfseries\color{GrisPizarra}\DocAutor} & &
  {\small\bfseries\color{GrisPizarra}\DocRevisor} \\[2pt]
  {\footnotesize\color{GrisMedio}\DocCargoAutor} & &
  {\footnotesize\color{GrisMedio}\DocCargoRevisor} \\[1pt]
  {\footnotesize\color{GrisMedio}\DocEntidad} & &
  {\footnotesize\color{GrisMedio}\DocEntidad} \\[1pt]
  {\footnotesize\color{GrisMedio}\DocFecha} & &
  {\footnotesize\color{GrisMedio}\DocFecha} \\
\end{tabular}

\end{document}
